
\subsubsection{Synthetic Test Data}

\textbf{CRITICAL DISCLAIMER}: This proof-of-concept uses synthetic test data for pipeline validation, NOT real GitHub commits. Results demonstrate technical functionality but have ZERO scientific validity for claims about real developer behavior.

Dataset specification:
\begin{itemize}
\item \textbf{Size}: 10,000 commits
\item \textbf{Dimensions}: 4 quality metrics
\item \textbf{Aspect Labels}: Balanced distribution (2,500 per aspect)
\item \textbf{Random Seed}: 42 (reproducible)
\item \textbf{Generation Method}: Randomized covariance structure without guaranteed aspect-dominant geometry
\end{itemize}

\textbf{Purpose}: Demonstrate that analysis pipeline executes correctly, gate evaluation logic functions properly, and permutation testing detects lack of signal when present.

\textbf{What This Validates}: (1) Methodology executes without errors, (2) Statistical tests work as designed, (3) Gate criteria are computable.

\textbf{What This Does NOT Validate}: Real developer behavior, architectural recommendations, field implications.

\subsubsection{Real Data Collection Protocol (Required for Publication)}

For scientific validity, the following is required:

\textbf{Phase 1A: Commit Mining (5 days)}
\begin{itemize}
\item Source: GitHub API (PyGitHub, PyDriller)
\item Target: 10,000 minimal-diff commits (AST distance < 20 nodes)
\item Filters: Python/JavaScript files, aspect-labeled commits, active repositories
\item Repository sampling: 500-1000 diverse repositories
\end{itemize}

\textbf{Phase 1B: Metric Computation (22 hours)}
\begin{itemize}
\item Δcorrectness: pytest/jest test pass rate changes
\item Δquality: SonarQube maintainability rating changes
\item Δsecurity: CodeQL alert count changes
\item Δefficiency: pytest-benchmark runtime changes
\item Infrastructure: SonarQube Docker instance, CodeQL CLI, 16-core CPU
\end{itemize}

\textbf{Phase 1C: Label Validation (2 days)}
\begin{itemize}
\item Expert annotation: n=500 commits with 3 annotators
\item Measure: Inter-rater agreement (Fleiss' kappa > 0.6), commit purity (> 70\% single-aspect)
\end{itemize}

\subsubsection{Research Questions}

\textbf{RQ1}: Do quality metrics exhibit statistical independence? (Test: coupling ≤ 0.2)

\textbf{RQ2}: Does covariance exhibit anisotropic aspect-dominant structure? (Test: gap > 2.0)

\textbf{RQ3}: Do aspect labels correspond to covariance structure? (Test: permutation p < 0.05)

\textbf{RQ4}: Are eigenvectors aligned with aspect-specific directions? (Test: z-score > 2.0)

\textbf{RQ5}: Is structure robust across repositories? (Test: LORO consistency ≥ 0.7)

\subsubsection{Computational Resources}

\begin{itemize}
\item Hardware: Single NVIDIA GPU (not used for analysis)
\item Software: Python 3.9, NumPy 1.24, SciPy 1.10
\item Runtime: Spectral analysis < 1 minute, Permutation test ~5 minutes
\item Storage: Synthetic dataset 80MB, Real dataset (future) ~500GB
\end{itemize}
